\setcounter{chapter}{1}
\chapter{Jordan測度}

\paragraph{本章の目標}
Jordan測度，Jordan可測集合，有限加法族（集合代数）

第1章では1次元のRiemann積分を導入し，Dirichlet関数によって
Riemann積分が常には定義できないことを見た．
その障害は，関数の不連続性だけではなく，
集合の「境界の大きさ」に由来している．
本章では，そのことをJordan測度の観点から整理する．

有界集合 $A \subset \RR^d$ の指示関数
\[
1_A(x) :=
\begin{cases}
1 & (x \in A),\\
0 & (x \notin A)
\end{cases}
\]
を考えると，$1_A$ がRiemann可積分であるかどうかは，
$A$ を有限個の直方体でどこまでうまく近似できるかという問題になる．
この「有限個」という条件が，
後のLebesgue理論との顕著な差である．

\section{\texorpdfstring{$\RR^d$}{Rd}の基本集合}

\begin{definition}[区間と長さ]
実数 $a \le b$ に対し，
\begin{align*}
[a,b] &:= \{x \in \RR \mid a \le x \le b\}, &
[a,b) &:= \{x \in \RR \mid a \le x < b\},\\
(a,b] &:= \{x \in \RR \mid a < x \le b\}, &
(a,b) &:= \{x \in \RR \mid a < x < b\}
\end{align*}
を $\RR$ の有界区間という．
これらの長さは端点の開閉に依らず
\[
|I| := b-a
\]
で定める．
\end{definition}

\begin{definition}[区間]
$\RR^d$ の有界区間とは，
各成分が有界区間である直積
\[
R = I_1 \times \cdots \times I_d
\]
のことである．
その体積を
\[
|R| := \prod_{k=1}^d |I_k|
\]
で定める．
$d=1,2,3$ のときは，それぞれ長さ，面積，体積と呼ぶ．
\end{definition}

\begin{remark}
端点の開閉は体積に影響しない．
以後，有限分割を扱う際には
\[
Q = \prod_{k=1}^d [a_k,b_k)
\]
という半開区間を主に用いる．
境界は後でJordan零集合として無視できるので，
閉区間・開区間・半開区間の違いは本質的でない．
\end{remark}

\begin{definition}[基本集合]
$\RR^d$ の基本集合とは，
有限個の互いに素な半開区間の合併
\[
E = \bigsqcup_{j=1}^n Q_j
\]
のことである．
このとき
\[
m_J(E) := \sum_{j=1}^n |Q_j|
\]
と定める．
\end{definition}

\begin{lemma}[格子分割]
有限個の有界区間の合併 $A = \sum_{i=1}^n B_i$は，基本集合$\sum_{j=1}^m Q_j$として書き直せる．
特に，書き換え後の体積$\sum_{j=1}^m |Q_j|$は分割の方法に依らない．
\end{lemma}
\begin{proof}
    （演習）
\end{proof}
%\begin{proof}
各 $B_i$ を
\[
B_i = I_{i,1} \times \cdots \times I_{i,d}
\]
と書く．
各座標方向 $k=1,\dots,d$ について，
$I_{i,k}$ に現れる全ての端点を集めて
\[
t_{k,0} < t_{k,1} < \cdots < t_{k,N_k}
\]
と並べる．
これらでできる格子の小半開区間
\[
Q_\alpha
:=
\prod_{k=1}^d [t_{k,\alpha_k-1},t_{k,\alpha_k})
\qquad
(1 \le \alpha_k \le N_k)
\]
を考える．

まず $B_i$ が半開区間である場合を考える．
各座標方向で
$[t_{k,\alpha_k-1},t_{k,\alpha_k})$ は
$I_{i,k}$ に含まれるか，さもなくば交わらないかのどちらかであるから，
各格子小区間 $Q_\alpha$ についても
\[
Q_\alpha \subset B_i
\quad\text{または}\quad
Q_\alpha \cap B_i = \emptyset
\]
のいずれかである．
したがって
\[
A
=
\bigsqcup_{Q_\alpha \subset A} Q_\alpha
\]
と表せるので，$A$ は基本集合として書ける．

一般の有界区間の場合も，端点の開閉の違いは格子の境界上でしか起こらないから，
同じ端点をもつ半開区間に直して上と同じ議論をすればよい．

次に体積和が表示に依らないことを示す．
$A$ の二つの基本集合表示
\[
A = \bigsqcup_{j=1}^m Q_j = \bigsqcup_{\ell=1}^r Q'_\ell
\]
を取る．
$Q_j$ と $Q'_\ell$ に現れる全ての端点を集めて共通の格子分割を作ると，
その各小区間 $S_\nu$ は各 $Q_j$ に対しても各 $Q'_\ell$ に対しても
「含まれる」か「交わらない」かのどちらかである．
ゆえに
\[
A = \bigsqcup_{S_\nu \subset A} S_\nu
\]
であり，しかも各 $Q_j$ も各 $Q'_\ell$ も
これらの $S_\nu$ の互いに素な合併に分解される．
したがって
\[
\sum_{j=1}^m |Q_j|
=
\sum_{S_\nu \subset A} |S_\nu|
=
\sum_{\ell=1}^r |Q'_\ell|
\]
となる．
よって，書き換え後の体積 $\sum_{j=1}^m |Q_j|$ は分割の方法に依らない．
\end{proof}


\begin{theorem}[基本集合の測度の性質]
基本集合 $E,F \subset \RR^d$ に対して次が成り立つ．
\begin{enumerate}
    \item （非負性）$m_J(E) \ge 0$，特に $m_J(\emptyset)=0$
    \item （単調性）$E \subset F$ ならば $m_J(E) \le m_J(F)$
    \item （加法性）$E \cap F = \emptyset$ ならば $m_J(E \sqcup F) = m_J(E) + m_J(F)$
    \item （平行移動不変性）任意の $c \in \RR^d$ に対して $m_J(E+c) = m_J(E)$
\end{enumerate}
\end{theorem}
\begin{proof}
補題より，$E,F$ に現れる全ての端点を集めた共通の格子分割
\[
\calP=\{R_1,\dots,R_N\}
\]
を取れば，ある添字集合 $I,J \subset \{1,\dots,N\}$ が存在して
\[
E=\bigsqcup_{i\in I} R_i,
\qquad
F=\bigsqcup_{i\in J} R_i
\]
と書ける．
このとき
\[
m_J(E)=\sum_{i\in I}|R_i|,
\qquad
m_J(F)=\sum_{i\in J}|R_i|
\]
である．

\begin{enumerate}
    \item 各 $|R_i| \ge 0$ であるから
    \[
    m_J(E)=\sum_{i\in I}|R_i| \ge 0.
    \]
    また $\emptyset=\bigsqcup_{i\in\emptyset}R_i$ だから
    \[
    m_J(\emptyset)=\sum_{i\in\emptyset}|R_i|=0.
    \]

    \item $E \subset F$ ならば，各 $R_i$ は互いに素であるから $I \subset J$ である．
    したがって
    \[
    m_J(F)-m_J(E)
    =\sum_{i\in J}|R_i|-\sum_{i\in I}|R_i|
    =\sum_{i\in J\setminus I}|R_i|
    \ge 0,
    \]
    よって $m_J(E)\le m_J(F)$ である．

    \item $E \cap F=\emptyset$ ならば $I\cap J=\emptyset$ であり，
    \[
    E\sqcup F=\bigsqcup_{i\in I\cup J}R_i
    \]
    だから
    \[
    m_J(E\sqcup F)
    =\sum_{i\in I\cup J}|R_i|
    =\sum_{i\in I}|R_i|+\sum_{i\in J}|R_i|
    =m_J(E)+m_J(F).
    \]

    \item $c=(c_1,\dots,c_d)\in\RR^d$ とする．
    各 $R_i=\prod_{k=1}^d[a_{i,k},b_{i,k})$ に対して
    \[
    R_i+c=\prod_{k=1}^d[a_{i,k}+c_k,b_{i,k}+c_k)
    \]
    であり，
    \[
    |R_i+c|
    =\prod_{k=1}^d\bigl((b_{i,k}+c_k)-(a_{i,k}+c_k)\bigr)
    =\prod_{k=1}^d(b_{i,k}-a_{i,k})
    =|R_i|.
    \]
    ゆえに
    \[
    E+c=\bigsqcup_{i\in I}(R_i+c)
    \]
    かつ
    \[
    m_J(E+c)
    =\sum_{i\in I}|R_i+c|
    =\sum_{i\in I}|R_i|
    =m_J(E).
    \]
\end{enumerate}
\end{proof}

\section{\texorpdfstring{$\RR^d$}{Rd}のRiemann積分}

第1章の1次元Riemann積分は，そのまま $d$ 次元へ拡張できる．
積分区間としては有界閉区間
\[
I = [a_1,b_1] \times \cdots \times [a_d,b_d]
\]
を取り，分割 $\calP$ は各座標方向の有限分割から得られる有限個の小区間の族
\[
I = \bigsqcup_{j=1}^N R_j
\]
とする．
分割の大きさは
\[
|\calP| := \max_{1 \le k \le d,\ 1 \le i \le n_k} (x_{k,i}-x_{k,i-1})
\]
で測る．

有界関数 $f : I \to \RR$ に対し，各小区間 $R_j$ から評価点 $t_j \in R_j$ を選んで
\[
S(f,\calP,t) := \sum_{j=1}^N f(t_j)\,|R_j|
\]
をRiemann和という．
またDarbouxの上和・下和を
\[
U(f,\calP) := \sum_{j=1}^N \sup_{x \in R_j} f(x)\,|R_j|,
\qquad
L(f,\calP) := \sum_{j=1}^N \inf_{x \in R_j} f(x)\,|R_j|
\]
で定める．
1次元の場合と同様に，Darboux上積分・下積分を
\[
\uint{I}{} f(x)\,dx := \inf_{\calP} U(f,\calP),
\qquad
\lint{I}{} f(x)\,dx := \sup_{\calP} L(f,\calP)
\]
で定める．
さらに
\[
\uint{I}{} f(x)\,dx = \lint{I}{} f(x)\,dx
\]
が成り立つとき，$f$ は $I$ 上Riemann可積分であるといい，
その共通値を
\[
\int_I f(x)\,dx
\]
と書く．
このことは，任意の $\eps > 0$ に対して
\[
U(f,\calP)-L(f,\calP)<\eps
\]
を満たす分割 $\calP$ が存在することと同値である．

特に $f = 1_A$ （指示関数）のとき，各小区間上での上限と下限は $0$ か $1$ しか取らないので，
\[
U(1_A,\calP)
= \sum_{R_j \cap A \neq \emptyset} |R_j|,
\qquad
L(1_A,\calP)
= \sum_{R_j \subset A} |R_j|
\]
となる．
つまり，指示関数の上和と下和は，
分割に沿って $A$ を外側・内側から近似する基本集合の体積そのものである．
これがJordan測度の出発点である．

\section{Jordan外測度・Jordan内測度}

\begin{definition}[Jordan外測度とJordan内測度]
集合 $A \subset \RR^d$ に対し，
\[
m_J^*(A)
:=
\inf\{m_J(E) \mid E \text{ は基本集合},\ A \subset E\}
\]
をJordan外測度，
\[
m_{J,*}(A)
:=
\sup\{m_J(E) \mid E \text{ は基本集合},\ E \subset A\}
\]
をJordan内測度という．
\end{definition}

\begin{remark}
Jordan理論は本質的に有界集合の理論である．
実際，$A$ が有界でなければ，有限個の区間で $A$ 全体を覆えないので
$m_J^*(A)=\infty$ となる．
後のLebesgue測度では，可算個の区間を使えるため，
この有界性の制約が大きく緩和される．
\end{remark}

\begin{proposition}
任意の集合 $A \subset \RR^d$ に対して
\[
m_{J,*}(A) \le m_J^*(A)
\]
が成り立つ．
\end{proposition}
\begin{proof}
$A$ を含む基本集合が存在しなければ $m_J^*(A)=\infty$ であり自明である．
そうでないとき，$E \subset A \subset F$ を満たす基本集合 $E,F$ に対して
基本集合の単調性より $m_J(E) \le m_J(F)$ である．
左辺で上限を，右辺で下限を取ればよい．
\end{proof}

\begin{definition}[Jordan可測，Jordan零集合]
有界集合 $A \subset \RR^d$ がJordan可測であるとは
\[
m_{J,*}(A) = m_J^*(A)
\]
が成り立つことをいう．
その共通値を $A$ のJordan測度と書き，$m_J(A)$ で表す．

また $m_J^*(N)=0$ を満たす集合 $N$ をJordan零集合という．
\end{definition}

\begin{remark}
厳密には，Darboux 上下積分は「固定した$I$の分割」による上和・下和で定義され，
Jordan 外測度・内測度は「任意の基本集合 $E,F$」で定義されているため，量化の範囲が異なる．
格子分割の補題を使うと，分割から基本集合が作れ，逆に基本集合は共通細分で分割の和として書けるので，両者が一致することが言える．
\end{remark}

\begin{lemma}
有界集合 $A \subset I$ に対して
\[
\uint{I}{} 1_A(x)\,dx = m_J^*(A),
\qquad
\lint{I}{} 1_A(x)\,dx = m_{J,*}(A)
\]
が成り立つ．
\end{lemma}
\begin{proof}
任意の分割 $\calP=\{R_1,\dots,R_N\}$ に対し
\[
G_{\calP}:=\bigcup_{R_j\cap A\neq\emptyset}R_j,
\qquad
F_{\calP}:=\bigcup_{R_j\subset A}R_j
\]
とおくと，$G_{\calP},F_{\calP}$ は基本集合で
\[
A\subset G_{\calP}\subset I,
\qquad
F_{\calP}\subset A
\]
を満たし，
\[
U(1_A,\calP)=m_J(G_{\calP}),
\qquad
L(1_A,\calP)=m_J(F_{\calP})
\]
である．
したがって任意の $\calP$ に対して
\[
m_J^*(A)\le U(1_A,\calP),
\qquad
L(1_A,\calP)\le m_{J,*}(A)
\]
となるから
\[
m_J^*(A)\le \uint{I}{}1_A(x)\,dx,
\qquad
\lint{I}{}1_A(x)\,dx\le m_{J,*}(A)
\]
を得る．

逆向きを示す．
まず $A\subset G$ を満たす基本集合 $G$ を取る．
$A\subset I$ だから $A\subset G\cap I$ であり，格子分割の補題により $G\cap I$ も基本集合である．
$G\cap I$ に現れる端点で $I$ を分割すると，各小区間は $G\cap I$ に含まれるか，さもなくば交わらないかのいずれかである．
この分割を $\calP$ とすると，$R_j\cap A\neq\emptyset$ ならば $R_j\subset G\cap I$ であるから
\[
U(1_A,\calP)
=m_J(G_{\calP})
\le m_J(G\cap I)
\le m_J(G).
\]
よって
\[
\uint{I}{}1_A(x)\,dx\le m_J(G)
\]
であり，右辺で $G$ に関する下限を取って
\[
\uint{I}{}1_A(x)\,dx\le m_J^*(A)
\]
を得る．

次に基本集合 $F\subset A$ を取る．
$F$ に現れる端点で $I$ を分割すると，$F$ はその分割の小区間の合併になる．
この分割を $\calP$ とすると，$F$ を構成する各小区間は $A$ に含まれるから
\[
L(1_A,\calP)\ge m_J(F).
\]
したがって
\[
\lint{I}{}1_A(x)\,dx\ge m_J(F)
\]
であり，右辺で $F$ に関する上限を取って
\[
\lint{I}{}1_A(x)\,dx\ge m_{J,*}(A)
\]
を得る．
以上より結論が従う．
\end{proof}

\begin{lemma}
基本集合の境界はJordan零集合である．
\end{lemma}
\begin{proof}
まず半開区間
\[
Q = \prod_{k=1}^d [a_k,b_k)
\]
について示す．
$\partial Q$ は有限個の座標超平面片の合併に含まれる：
\[
\partial Q \subset
\bigcup_{k=1}^d
\Bigl(
\{x_k=a_k\} \cap \prod_{i=1}^d [a_i,b_i]
\Bigr)
\cup
\bigcup_{k=1}^d
\Bigl(
\{x_k=b_k\} \cap \prod_{i=1}^d [a_i,b_i]
\Bigr).
\]
各面は，厚さ $\delta$ の薄い板で覆えるので，
その被覆の体積は $\delta$ に比例していくらでも小さくできる．
したがって各面はJordan零であり，有限和である $\partial Q$ もJordan零である．
基本集合の境界は有限個の半開区間の境界の和集合に含まれるから，同様にJordan零である．
\end{proof}

\begin{theorem}[Jordan可測の判定法]
有界集合 $A \subset \RR^d$ と，それを含む有界閉区間 $I$ を取る．
このとき次は同値である．
\begin{enumerate}
    \item $1_A$ は $I$ 上でRiemann可積分である
    \item $A$ はJordan可測である
    \item 任意の $\eps > 0$ に対し，基本集合 $F,G$ が存在して
    \[
    F \subset A \subset G,
    \qquad
    m_J(G \setminus F) < \eps
    \]
    を満たす
    \item 境界 $\partial A := \cl A \setminus \intr A$ はJordan零集合である
\end{enumerate}
\end{theorem}

\begin{remark}
閉包 $\cl A$ とは $A$ を含む最小の閉集合 $\bigcap \{ F \mid A \subset F, F \text { closed} \}$ であり，
内部 $\intr A$ とは $A$ に含まれる最大の開集合 $\bigcup \{ U \mid U \subset A, \text{ open} \}$である．
同値な言い方をすれば，
$x \in \cl A$ とは $x$ の任意の近傍が $A$ と交わること，
$x \in \intr A$ とは $x$ のある近傍が $A$ に含まれることをいう．
\end{remark}

\begin{proof}
\textbf{(1) $\Leftrightarrow$ (2)}\quad
直前の補題より
\[
\uint{I}{}1_A(x)\,dx = m_J^*(A),
\qquad
\lint{I}{}1_A(x)\,dx = m_{J,*}(A)
\]
である．
したがって
\[
1_A \text{ が } I \text{ 上Riemann可積分}
\]
であること，すなわち
\[
\uint{I}{}1_A(x)\,dx = \lint{I}{}1_A(x)\,dx
\]
であることは，
\[
m_J^*(A)=m_{J,*}(A)
\]
すなわち $A$ がJordan可測であることと同値である．

\textbf{(1) $\Rightarrow$ (3)}\quad
$1_A$ がRiemann可積分なら，Darbouxの判定法により，
任意の $\eps > 0$ に対して
\[
U(1_A,\calP) - L(1_A,\calP) < \eps
\]
となる分割 $\calP$ が取れる．
この分割に対し
\[
F := \bigcup_{R_j \subset A} R_j,
\qquad
G := \bigcup_{R_j \cap A \neq \emptyset} R_j
\]
とおけば，$F,G$ は基本集合で
$F \subset A \subset G$ かつ
\[
m_J(G \setminus F)
=
U(1_A,\calP) - L(1_A,\calP)
< \eps
\]
である．

\textbf{(3) $\Rightarrow$ (2)}\quad
任意の $\eps > 0$ に対して
\[
F \subset A \subset G,
\qquad
m_J(G \setminus F) < \eps
\]
を満たす基本集合 $F,G$ が存在するとする．
このとき
\[
m_{J,*}(A)\ge m_J(F),
\qquad
m_J^*(A)\le m_J(G)
\]
だから
\[
0 \le m_J^*(A)-m_{J,*}(A)
\le m_J(G)-m_J(F)
= m_J(G \setminus F)
< \eps.
\]
$\eps$ の任意性より $m_J^*(A)=m_{J,*}(A)$ である．

\textbf{(3) $\Rightarrow$ (4)}\quad
$F \subset A \subset G$ とすると
\[
\partial A \subset \cl G \setminus \intr F
\subset (G \setminus F) \cup \partial G \cup \partial F.
\]
ここで $\partial F,\partial G$ はJordan零集合であるから，
\[
m_J^*(\partial A)
\le m_J(G \setminus F) + m_J^*(\partial G) + m_J^*(\partial F)
= m_J(G \setminus F)
< \eps
\]
が従う．
したがって $\partial A$ はJordan零である．

\textbf{(4) $\Rightarrow$ (3)}\quad
$\partial A$ がJordan零集合なので，
任意の $\eps > 0$ に対して
\[
\partial A \subset H,
\qquad
H \text{ は基本集合},
\qquad
m_J(H) < \eps
\]
となる基本集合 $H$ が存在する．
いま
\[
K := \cl A \setminus H
\]
とおくと，$\cl A \setminus \partial A = \intr A$ より $K \subset \intr A$ である．
$K$ はコンパクトで $\intr A$ は開集合だから，
$K$ を有限個の小半開区間で覆って，その合併 $F$ を
\[
K \subset F \subset \intr A \subset A
\]
となるように取れる．
すると $A \subset F \cup H$ なので，
$G := F \cup H$ とおけば基本集合で
\[
F \subset A \subset G,
\qquad
G \setminus F \subset H
\]
が成り立つ．
したがって
\[
m_J(G \setminus F) \le m_J(H) < \eps.
\]
\end{proof}

\begin{corollary}
有界集合 $A \subset I$ がJordan可測であるとき，
\[
\int_I 1_A(x)\,dx = m_J(A)
\]
が成り立つ．
\end{corollary}
\begin{proof}
上の定理より $1_A$ はRiemann可積分である．
さらに基本集合 $F,G$ による近似に対して
\[
m_J(F) \le \int_I 1_A(x)\,dx \le m_J(G)
\]
が成り立つので，$m_J(G \setminus F)$ を $0$ に近づければ
積分値は $m_J(A)$ に一致する．
\end{proof}

\begin{example}
\begin{enumerate}
    \item 基本集合はJordan可測であり，そのJordan測度は最初に定めた $m_J$ に一致する
    \item 有限集合はJordan零集合である
    \item 線分や超平面片など，$\RR^d$ の中で厚みをもたない典型的な図形はJordan零集合である
\end{enumerate}
\end{example}
\begin{proof}
\begin{enumerate}
    \item 基本集合 $E$ 自身を内側近似と外側近似に取れば
    $m_{J,*}(E)=m_J^*(E)=m_J(E)$ である
    \item 1点集合は十分小さい区間1つで覆えるのでJordan零である．有限集合はその有限和である
    \item 例えば超平面片 $\{x_d=0\} \cap [-M,M]^d$ は厚さ $\delta$ の板
    \[
    [-M,M]^{d-1} \times [-\delta,\delta]
    \]
    で覆えるのでJordan零である．線分も同様に細い区間で覆える
\end{enumerate}
\end{proof}

\begin{remark}
差が最も見えやすい例は
\[
A = [0,1] \cap \QQ
\]
である．
各点集合 $\{q\}$ はJordan零集合だが，
$A$ は $[0,1]$ の中で稠密なので
\[
\partial A = [0,1]
\]
となり，Jordan可測ではない．
したがってJordan可測集合は可算和で閉じていない．

一方で，後に導入するLebesgue測度では
$A$ は可測で，しかも測度 $0$ をもつ．
ここに「有限分割に基づく理論」と「可算分割に耐える理論」の違いが現れる．
\end{remark}

\begin{remark}
上の例は「有界集合でもJordan可測とは限らない」ことを示すが，
この集合は閉集合ではない．
閉集合や開集合が必ずJordan可測だと言っているわけでもないことに注意する．
Jordan可測性の本質は，開閉ではなく
「境界がJordan零かどうか」にある．
\end{remark}

\section{Jordan測度の性質}

\begin{theorem}[Jordan可測集合の集合代数的閉性]
Jordan可測集合 $A,B \subset \RR^d$ と $c \in \RR^d$ に対して，
\[
A \cup B,\qquad
A \cap B,\qquad
A \setminus B,\qquad
A+c
\]
は再びJordan可測である．
\end{theorem}
\begin{proof}
境界について
\[
\partial(A \cup B) \subset \partial A \cup \partial B,
\qquad
\partial(A \cap B) \subset \partial A \cup \partial B,
\]
\[
\partial(A \setminus B) \subset \partial A \cup \partial B,
\qquad
\partial(A+c) = \partial A + c
\]
が成り立つ．
Jordan零集合の有限和と平行移動は再びJordan零集合だから，
判定法より各集合はJordan可測である．
\end{proof}

\begin{theorem}[Jordan測度の性質]
Jordan可測集合 $A,B \subset \RR^d$ と $c \in \RR^d$ に対して次が成り立つ．
\begin{enumerate}
    \item $m_J(A) \ge 0$，特に $m_J(\emptyset)=0$
    \item $A \subset B$ ならば $m_J(A) \le m_J(B)$
    \item $A \cap B = \emptyset$ ならば
    \[
    m_J(A \sqcup B) = m_J(A) + m_J(B)
    \]
    が成り立つ
    \item 一般に
    \[
    m_J(A \cup B) = m_J(A) + m_J(B) - m_J(A \cap B)
    \]
    が成り立つ
    \item $m_J(A+c)=m_J(A)$
    \item $m_J(\cl A)=m_J(A)=m_J(\intr A)$
\end{enumerate}
\end{theorem}
\begin{proof}
\begin{enumerate}
    \item 外測度の定義から明らかである．また $\emptyset$ は基本集合なので測度は $0$ である
    \item $A \subset B$ ならば，$B$ を覆う基本集合はすべて $A$ も覆うので
    $m_J^*(A) \le m_J^*(B)$ である．
    両者はJordan可測だから
    \[
    m_J(A)=m_J^*(A) \le m_J^*(B)=m_J(B)
    \]
    となる
    \item ある有界閉区間 $I$ が $A \cup B$ を含むように取る．
    すると
    \[
    1_{A \sqcup B} = 1_A + 1_B
    \]
    であり，左辺右辺はいずれも $I$ 上Riemann可積分である．
    よって第1章の積分の線形性から
    \[
    m_J(A \sqcup B)
    = \int_I 1_{A \sqcup B}(x)\,dx
    = \int_I 1_A(x)\,dx + \int_I 1_B(x)\,dx
    = m_J(A)+m_J(B)
    \]
    を得る
    \item 恒等式
    \[
    1_{A \cup B} = 1_A + 1_B - 1_{A \cap B}
    \]
    を積分すればよい
    \item 基本集合の体積は平行移動不変なので，$m_J^*(A+c)=m_J^*(A)$，
    $m_{J,*}(A+c)=m_{J,*}(A)$ が成り立つ．
    よってJordan測度も不変である
    \item
    \[
    \intr A \subset A \subset \cl A,
    \qquad
    \cl A \setminus \intr A = \partial A
    \]
    であり，$\partial A$ はJordan零集合である．
    したがって (2), (3) より
    \[
    m_J(\cl A) = m_J(A) = m_J(\intr A)
    \]
    が従う
\end{enumerate}
\end{proof}

\begin{remark}
Jordan測度は，Jordan可測集合の上では
「有限加法的な体積」として振る舞う．
しかし，その本質はあくまで有限加法性であり，
可算個の集合を同時に扱うところで理論が止まる．
ここが後のLebesgue測度との最大の違いである．
\end{remark}

\begin{example}
\begin{enumerate}
    \item $\QQ \cap [0,1]$ を
    \[
    \QQ \cap [0,1] = \bigcup_{n=1}^\infty \{q_n\}
    \]
    と可算個の1点集合の和に書く．
    各 $\{q_n\}$ はJordan零集合，したがってJordan可測であるが，
    和集合 $\QQ \cap [0,1]$ はJordan可測でない．
    実際，$\QQ \cap [0,1]$ もその補集合 $([0,1]\setminus\QQ)$ も $[0,1]$ で稠密だから，
    \[
    \partial(\QQ \cap [0,1])=[0,1]
    \]
    となる．
    したがってJordan可測集合の全体は可算和で閉じていない．

    \item 上と同じ列 $q_1,q_2,\dots$ を用いて
    \[
    E_n:=\{q_1,\dots,q_n\}
    \]
    とおくと，$E_n$ はJordan可測で
    \[
    m_J(E_n)=0 \qquad (n=1,2,\dots)
    \]
    であり，しかも
    \[
    E_1 \subset E_2 \subset \cdots,
    \qquad
    \bigcup_{n=1}^\infty E_n=\QQ \cap [0,1]
    \]
    である．
    もしLebesgue測度のような可算加法性や下からの連続性がJordan測度にも成り立つなら，
    \[
    m_J\Bigl(\bigcup_{n=1}^\infty E_n\Bigr)
    = \sum_{n=1}^\infty m_J(E_n\setminus E_{n-1})
    = 0
    \]
    あるいは
    \[
    m_J\Bigl(\bigcup_{n=1}^\infty E_n\Bigr)
    = \lim_{n\to\infty} m_J(E_n)=0
    \]
    と書きたくなる．
    しかし左辺の集合 $\QQ \cap [0,1]$ 自体がJordan可測でないので，
    これらの等式はそもそも意味をもたない．
    つまりJordan理論では，可算極限に進んだところで理論そのものが止まる．
\end{enumerate}
\end{example}

\section{有限加法族}

\begin{definition}[有限加法族]
集合 $X$ の部分集合族 $\calA \subset \mathcal P(X)$ が有限加法族であるとは，
次の3条件を満たすことをいう．
\begin{enumerate}
    \item $\emptyset \in \calA$
    \item $A \in \calA$ ならば $X \setminus A \in \calA$
    \item $A,B \in \calA$ ならば $A \cup B \in \calA$
\end{enumerate}
有限加法族は集合代数（algebra of sets）とも呼ばれる．
\end{definition}

\begin{proposition}[有限加法族の基本性質]
$\calA$ が $X$ 上の有限加法族ならば，任意の $A,B,A_1,\dots,A_n \in \calA$ に対して
\begin{enumerate}
    \item $X \in \calA$
    \item $A \cap B \in \calA$
    \item $A \setminus B \in \calA$
    \item $\bigcup_{k=1}^n A_k \in \calA$
    \item $\bigcap_{k=1}^n A_k \in \calA$
\end{enumerate}
が成り立つ．
\end{proposition}
\begin{proof}
\[
X = X \setminus \emptyset,
\qquad
A \cap B = X \setminus \bigl((X \setminus A) \cup (X \setminus B)\bigr),
\qquad
A \setminus B = A \cap (X \setminus B)
\]
だから，(1)--(3) は補集合と和集合についての閉性から従う．
(4), (5) はこれを有限回繰り返せばよい．
\end{proof}

\begin{definition}[有限加法的集合関数]
有限加法族 $\calA$ 上の集合関数 $\mu : \calA \to [0,\infty)$ が有限加法的であるとは，
$A,B \in \calA$ が互いに素であるとき
\[
\mu(A \sqcup B) = \mu(A)+\mu(B)
\]
が成り立つことをいう．
\end{definition}

\begin{theorem}[Jordan可測集合の有限加法族]
有界Jordan可測集合 $X \subset \RR^d$ を固定し，
\[
\calA_X := \{A \subset X \mid A \text{ はJordan可測}\}
\]
とおく．
このとき $\calA_X$ は $X$ 上の有限加法族であり，
$m_J : \calA_X \to [0,\infty)$ は有限加法的である．
\end{theorem}
\begin{proof}
$\emptyset$ はJordan可測である．
また $A \subset X$ がJordan可測なら
\[
\partial(X \setminus A) \subset \partial X \cup \partial A
\]
だから $X \setminus A$ もJordan可測である．
さらに $A,B \subset X$ がJordan可測なら
\[
\partial(A \cup B) \subset \partial A \cup \partial B
\]
より $A \cup B$ もJordan可測である．
したがって $\calA_X$ は有限加法族である．
有限加法性は前節のJordan測度の有限加法性そのものである．
\end{proof}

\begin{remark}
注意すべきなのは，
Jordan可測集合全体をそのまま $\RR^d$ 上で見ても有限加法族にはならないことである．
補集合を取ると一般に有界性を失うからである．
したがって，Jordan理論では最初に有界な全体集合 $X$ を固定して
その内部で議論するのが自然である．
\end{remark}

\begin{remark}
さらに $\calA_X$ は一般には $\sigma$-加法族ではない．
例えば $X=[0,1]$ とし，$\QQ \cap [0,1]=\{q_1,q_2,\dots\}$ と書くと，
各単点集合 $\{q_n\}$ は $\calA_X$ に属し，
\[
m_J(\{q_n\})=0
\]
であるが，
\[
\bigcup_{n=1}^\infty \{q_n\} = \QQ \cap [0,1]
\]
はJordan可測でない．
この欠点を克服するために，
次章以降では「有限加法族」を「$\sigma$-加法族」へ，
Jordan測度をLebesgue測度へと拡張していく．
\end{remark}
